% ============================================================
% §04 — Domain Kit Model
% Author: Ken (Software Architect)
% Integrated by: Leslie (LaTeX Engineer)
% Last updated: 2026-04-20
% ============================================================

\chapter{Domain Kit Model}
\label{ch:domain-kits}

% ============================================================
\section{Purpose and Scope}
\label{sec:kit-purpose}
% ============================================================

The gert v2 core runtime is deliberately domain-agnostic.  It executes steps,
enforces governance, captures evidence, and writes traces---but it carries no
opinion about what kind of work a runbook describes.  This is a feature, not an
omission: a domain-agnostic kernel is the prerequisite for a stable, long-lived
runtime contract.

Domain Kits are the mechanism by which domain-specific semantics enter the
system without contaminating the kernel.  A Domain Kit is a higher-level
semantic layer built atop core runtime primitives.  It provides five things:

\begin{enumerate}
  \item \textbf{Authoring schemas.}  A Kit defines a domain-specific YAML
    vocabulary---specialised step types, field structures, validation
    constraints---that is natural for authors working in that domain.  An
    incident-response Kit might offer a \texttt{triage} step type with severity
    and escalation fields; an operations Kit might offer a
    \texttt{change-request} type with approval-matrix fields.

  \item \textbf{Validation rules.}  A Kit contributes domain-specific semantic
    validators that run at parse time or plan time.  These validators enforce
    invariants that the core schema cannot express: ``a rollback step must
    follow every deploy step,'' ``every incident runbook must declare a DRI,''
    ``escalation timeout must not exceed SLA.''

  \item \textbf{Lowering (compilation into core primitives).}  Kit-specific
    step types do not reach the runtime.  A Kit compiler transforms (lowers)
    its domain-specific AST nodes into sequences of core step types
    (\texttt{cli}, \texttt{tool}, \texttt{manual}, \texttt{branch},
    \texttt{iterate}) before the planner produces an ExecutionPlan.  The
    runtime executes only core primitives; it is never aware that a Kit was
    involved.

  \item \textbf{Projections and read models.}  A Kit may define read-only
    projections over run traces.  An operations Kit might project a ``change
    log'' view from the raw trace events; an incident Kit might project a
    ``timeline'' view.  Projections are computed after execution; they do not
    affect runtime behaviour.

  \item \textbf{Testing contracts.}  A Kit defines what ``correct'' means for
    its domain: fixture runbooks, expected lowered output, expected validation
    errors, expected projection shapes.  These contracts are the Kit's
    acceptance test suite.
\end{enumerate}

Domain Kits do not contribute runtime capabilities.  They do not register
tools, subscribe to events, or modify execution behaviour.  Those are extension
concerns (\S\ref{ch:extension}).  A Kit operates entirely at the authoring and
analysis layer---it shapes what authors write and what reviewers see, not what
the engine does.

% ============================================================
\section{Architectural Narrative}
\label{sec:kit-narrative}
% ============================================================

The v2 architecture draws a hard boundary between three concerns: what the
engine executes (core runtime), what capabilities are available at runtime
(extensions and tools), and what vocabulary authors use to express intent
(Domain Kits).  The first two are well-defined in the existing design; the
third has been implicit until now.

Without Domain Kits, domain semantics have two bad places to live.  They can
leak into the core schema, which makes the kernel unstable---every new domain
(incident response, change management, compliance audits, database operations)
adds fields and step types to the core, bloating it into an ever-expanding
union type.  Or they can be pushed into extensions, which conflates authoring
concerns with runtime concerns---an extension that contributes domain-specific
validators, authoring schemas, and projections is doing fundamentally different
work than an extension that contributes tools or event subscriptions.

Domain Kits solve this by providing a clean, well-defined home for domain
semantics.  The design principle is:

\begin{description}
  \item[Core owns execution.] The core runtime
    (\S\ref{ch:architecture}), governance layer (\S\ref{ch:governance}),
    evidence/tracing system (\S\ref{ch:evidence}), and event model
    (\S\ref{ch:events}) are stable, domain-agnostic infrastructure.

  \item[Kits own vocabulary.] Authoring schemas, domain validation rules,
    lowering/compilation, projections, and testing contracts belong to Kits.  A
    Kit is the right place to encode ``what does a well-formed incident runbook
    look like?'' or ``what fields does a change-request step require?''

  \item[Extensions own capabilities.] Tool registration, event subscription,
    schema extension (namespaced \texttt{x-*} fields), policy contribution, and
    provider registration belong to extensions (\S\ref{ch:extension}).

  \item[Tools own effects.] Executable actions---invoking binaries, calling
    APIs, performing side-effecting work---belong to the tool runtime
    (\S\ref{ch:tools}).
\end{description}

This separation preserves the kernel's stability guarantee: the core schema and
runtime contract do not change when a new domain is added.  The Kit compiles
domain-specific vocabulary down to core primitives; the runtime never sees
Kit-specific nodes.

% ============================================================
\section{The Clean Kernel Principle}
\label{sec:kit-clean-kernel}
% ============================================================

The key architectural invariant is: \textbf{the runtime never executes
Kit-specific step types.}  All domain-specific authoring nodes are lowered
(compiled) into core primitives before the planner sees them.

The compilation pipeline is:

\begin{verbatim}
Kit YAML → Kit Parser → Kit AST → Kit Compiler → Core YAML → Core Parser → ExecutionPlan
\end{verbatim}

Or equivalently, when integrated into the gert pipeline:

\begin{enumerate}
  \item The user authors a runbook using Kit vocabulary (e.g.,
    \texttt{apiVersion: runbook/v2+ops/v1}).
  \item The Kit's parser validates domain-specific fields and produces a
    Kit-level AST.
  \item The Kit's compiler lowers the AST into a standard \texttt{runbook/v2}
    document using only core step types.
  \item The core parser and planner process the lowered document normally.
  \item The runtime executes the plan with no knowledge that a Kit was
    involved.
\end{enumerate}

This is a compile-time transformation, not a runtime interception.  The lowered
output is a valid \texttt{runbook/v2} document that can be inspected,
version-controlled, and executed without the Kit installed.  This property is
critical for portability: a Kit-authored runbook can always be reduced to its
core equivalent.

The Kit compiler must be deterministic: the same Kit input must always produce
the same core output.  This ensures that \texttt{gert test} scenarios written
against lowered output remain stable across Kit versions.

% ============================================================
\section{Layer Relationships}
\label{sec:kit-layers}
% ============================================================

The v2 system has four distinct layers.  Each layer has a defined
responsibility boundary and a defined relationship to the others.

\begin{table}[ht]
\centering
\caption{Four-layer responsibility matrix}
\label{tab:kit-layers}
\small
\begin{tabular}{p{0.15\linewidth}p{0.30\linewidth}p{0.25\linewidth}p{0.18\linewidth}}
\toprule
\textbf{Layer} & \textbf{Owns} & \textbf{Does NOT own} & \textbf{Runs at} \\
\midrule
Core Runtime &
  Execution, state machine, governance enforcement, evidence, event emission, trace writing &
  Domain vocabulary, tool implementations, UI rendering &
  Run time \\
\midrule
Domain Kits &
  Authoring schemas, domain validation, lowering/compilation, projections, testing contracts &
  Execution, runtime capabilities, side effects &
  Author time / compile time \\
\midrule
Extension Runtime &
  Runtime capabilities: tool registration, event subscription, schema extension (\texttt{x-*}), policy contribution, provider registration &
  Authoring vocabulary, domain semantics, execution control &
  Run time (out-of-process) \\
\midrule
Tool Runtime &
  Side-effecting actions: binary execution, API calls, MCP tool invocation &
  Governance, evidence, schema, authoring vocabulary &
  Run time (per-invocation or persistent process) \\
\bottomrule
\end{tabular}
\end{table}

Key distinctions:

\begin{description}
  \item[Domain Kits are NOT extensions.] Extensions contribute runtime
    capabilities via the capability grant model
    (\S\ref{sec:ext-capabilities}).  Kits contribute authoring semantics via
    compilation.  An extension runs as a child process during execution; a Kit
    runs as a compiler before execution.  They serve fundamentally different
    purposes and operate at different phases.

  \item[Domain Kits are NOT tools.] Tools perform side-effecting work.  Kits
    are pure transformations with no side effects.  A Kit compiler reads YAML
    and writes YAML; it never invokes binaries, calls APIs, or modifies system
    state.

  \item[Extensions may accompany a Kit but are distinct.] A Kit distribution
    (e.g., an ``incident response pack'') may ship both a Domain Kit (authoring
    schemas, lowering rules) and an extension (incident-specific tools,
    PagerDuty providers).  These are bundled for convenience but are
    architecturally separate: the Kit operates at compile time, the extension
    operates at run time.
\end{description}

% ============================================================
\section{Kit Manifest and Compatibility}
\label{sec:kit-manifest}
% ============================================================

A Domain Kit ships a \texttt{gert-kit.yaml} manifest alongside its compiler.
The manifest declares identity, compatibility, and the Kit's schema
contributions.

\begin{minted}{yaml}
# gert-kit.yaml — Domain Kit Manifest
apiVersion: kit/v1

meta:
  name: gert.ops                     # Reverse-DNS namespaced identifier
  version: "1.0.0"                   # Semver
  description: "Operations runbook authoring kit"
  author: "Gert Project <gert@ormasoft.dev>"

compatibility:
  gert-core: ">=2.0.0 <3.0.0"       # Minimum gert core version (semver range)

schema:
  apiVersion-suffix: "ops/v1"        # Runbooks authored with this Kit use
                                     # apiVersion: runbook/v2+ops/v1
  json-schema: "./schemas/ops-v1.json"  # Kit-specific JSON Schema overlay

compiler:
  executable: "./bin/gert-kit-ops"   # Kit compiler binary
  args: ["compile"]
  input-format: yaml                 # Kit YAML in, core YAML out
  output-format: yaml

step-types:                          # Domain-specific step types this Kit defines
  - name: change-request
    lowers-to: [manual, tool]        # Core types this compiles into
  - name: rollback
    lowers-to: [cli, branch]
  - name: triage
    lowers-to: [manual, branch]

validators:                          # Domain-specific validation rules
  - name: ops/rollback-follows-deploy
    phase: plan                      # parse | plan
    description: "Every deploy step must be followed by a rollback step"
  - name: ops/dri-required
    phase: parse
    description: "Every ops runbook must declare a DRI in meta.roles"

projections:                         # Read models over trace data
  - name: ops/change-log
    description: "Change management timeline derived from trace events"
  - name: ops/evidence-summary
    description: "Audit-ready evidence summary grouped by governance primitive"
\end{minted}

\paragraph{Compatibility and versioning rules.}

\begin{itemize}
  \item A Kit declares a semver range for the minimum gert core version in
    \texttt{compatibility.gert-core}.  If the host version falls outside this
    range, the Kit is rejected at compile time with a structured error.

  \item Breaking changes to a Kit's authoring schema require a Kit major
    version bump.  The \texttt{apiVersion-suffix} encodes the Kit's major
    version (e.g., \texttt{ops/v1} $\to$ \texttt{ops/v2}).  This is
    independent of the gert core version.

  \item Additive changes (new optional fields, new step types) within a Kit
    major version are backward-compatible and do not require a version bump to
    the suffix.

  \item The Kit's \texttt{json-schema} artifact carries a semver artifact
    version for minor/patch tracking, following the same pattern as core schema
    artifacts (\S\ref{subsec:versioning-strategy}).
\end{itemize}

\paragraph{Local-first constraints.}

Domain Kits are portable semantic layers.  They run on any supported gert
host---a local laptop, a CI runner, a self-hosted server---without requiring
network access, cloud services, or central registries.  A Kit is a directory
containing a manifest, a compiler binary, and schema artifacts.  It can be
vendored into a repository, distributed as a tarball, or installed from a
package manager.  There is no Kit registry service, no Kit telemetry, and no
phone-home behaviour.  This aligns with gert's local-first design philosophy
(\S\ref{ch:overview}).

% ============================================================
\section{Extension Runtime Audit}
\label{sec:kit-ext-audit}
% ============================================================

The following table audits concerns currently described in the extension
runtime chapter (\S\ref{ch:extension}) and evaluates whether each is properly
an extension concern, a Domain Kit concern, or shared.

\begin{longtable}{p{0.15\linewidth}p{0.13\linewidth}p{0.14\linewidth}p{0.46\linewidth}}
\toprule
\textbf{Concern} & \textbf{Current location} & \textbf{Recommendation} & \textbf{Rationale} \\
\midrule
\endfirsthead

\toprule
\textbf{Concern} & \textbf{Current location} & \textbf{Recommendation} & \textbf{Rationale} \\
\midrule
\endhead

\midrule
\multicolumn{4}{r}{\footnotesize\itshape Continued on next page} \\
\bottomrule
\endfoot

\bottomrule
\endlastfoot

Domain-specific validators &
  Implicit in \texttt{capability/\allowbreak schema-extension} &
  \textbf{Move to Kit.} &
  Validation rules that enforce domain invariants (e.g., ``rollback must follow
  deploy'') are authoring-time concerns.  They belong in the Kit's validation
  phase, not in a runtime extension.  The extension
  \texttt{capability/\allowbreak schema-extension} should be reserved for namespaced
  field contributions (\texttt{x-*}), not for semantic validation of domain
  workflows. \\

Authoring schema registration &
  \texttt{capability/\allowbreak schema-extension} &
  \textbf{Clarify boundary.} &
  Extensions contribute namespaced runtime-observable fields (\texttt{x-acme.*}).
  Kits contribute domain-specific authoring vocabulary (new step types, required
  field structures).  These are different operations.  Add a clarifying note to
  \S\ref{sec:ext-capabilities} that \texttt{capability/\allowbreak schema-extension} is
  for runtime-observable namespace extensions, not for authoring-vocabulary
  contributions. \\

Projections / read models &
  Not currently addressed &
  \textbf{Kit concern.} &
  Projections are pure, read-only transformations over trace data.  They are
  computed after execution and do not require runtime capabilities.  They belong
  in the Kit layer.  Extensions that need to observe execution should use
  \texttt{capability/\allowbreak event-subscription}; post-hoc analysis belongs to
  Kits. \\

Compilation hooks &
  Not currently addressed &
  \textbf{Kit concern.} &
  The lowering/compilation step is a pure transformation that happens before the
  runtime starts.  It has no runtime footprint and requires no capabilities.
  Compilation is entirely within the Kit layer. \\

Domain testing fixtures &
  Not currently addressed &
  \textbf{Kit concern.} &
  Test fixtures for domain-specific authoring (example runbooks, expected lowered
  output, expected validation errors) are Kit deliverables.  They do not interact
  with the extension runtime. \\

Tool registration &
  \texttt{capability/\allowbreak tool-registration} &
  \textbf{Stays in extension runtime.} &
  Tool registration is a runtime capability.  A Kit distribution may ship an
  accompanying extension that registers domain-specific tools, but tool
  registration itself remains an extension concern. \\

Policy contribution &
  \texttt{capability/\allowbreak policy-contribution} &
  \textbf{Stays in extension runtime.} &
  Policy rules that are evaluated at run time (pre-step, pre-run) are runtime
  concerns.  Domain-specific validation rules that are evaluated at parse/plan
  time belong to Kits; runtime policy enforcement belongs to extensions. \\

Event subscription &
  \texttt{capability/\allowbreak event-subscription} &
  \textbf{Stays in extension runtime.} &
  Event observation is inherently a runtime activity.  No change needed. \\

\end{longtable}

\paragraph{Summary of changes needed in existing sections.}

\begin{enumerate}
  \item Add a clarifying paragraph to \texttt{capability/\allowbreak schema-extension}
    (\S\ref{sec:ext-capabilities}) noting the boundary: extensions contribute
    namespaced runtime fields; Kits contribute authoring vocabulary and domain
    step types.

  \item Add a forward reference from the extension runtime chapter to the
    Domain Kit chapter for readers who arrive at the extension spec looking for
    domain-specific validation or projections.

  \item No capabilities need to be moved or removed from the extension
    taxonomy.  The taxonomy is correct for runtime concerns; the issue was that
    domain/authoring concerns had no defined home.  Domain Kits provide that
    home.
\end{enumerate}

% ============================================================
\section{Kit-Zero: The DRI Operations Model}
\label{sec:kit-zero}
% ============================================================

The current DRI\,/\,operations runbook model---the step types, governance
primitives, evidence capture patterns, and approval workflows defined in the
core schema (\S\ref{ch:schema}) and governance chapter
(\S\ref{ch:governance})---is gert's first Domain Kit.  It is Kit-0: the
operations domain that motivated gert's existence.

In the current design, Kit-0 is implicit.  Its step types (\texttt{cli},
\texttt{manual}, \texttt{tool}, \texttt{branch}, \texttt{iterate}), governance
primitives (command allowlists, env blocking, redaction, approval gates), and
evidence types (text, choice, attachment with SHA256) are baked into the core
schema.  This is acceptable for v2.0 because the operations domain is gert's
primary use case and the Kit model is being introduced as a roadmap item, not a
v2.0 deliverable.

The architectural intent is clear: as the Kit model matures (v2.1+), the
operations-specific authoring vocabulary should be extractable into a
standalone Kit (\texttt{gert.ops}) that compiles into core primitives.  The
core schema would then contain only the minimal execution primitives (step
sequencing, variable binding, governance hooks, evidence slots), and the
operations vocabulary (DRI roles, change-request workflows, incident triage
patterns) would live in the Kit.  This extraction is not required for v2.0 but
should be planned for.

Kit-0 grounds the Domain Kit concept in something concrete.  When explaining
Kits to implementors or users, the answer to ``what does a Kit look like?'' is:
``look at the operations runbook model you already use---that is a Kit.''

% ============================================================
\section{Non-Goals}
\label{sec:kit-nongoals}
% ============================================================

Domain Kits are a precisely scoped mechanism.  The following are explicitly not
in scope:

\begin{description}
  \item[Kits are not extensions.] A Kit does not contribute runtime
    capabilities.  It does not register tools, subscribe to events, modify
    execution behaviour, or run as a child process during execution.  The
    extension capability taxonomy (\S\ref{sec:ext-capabilities}) does not
    apply to Kits.

  \item[Kits are not tools.] A Kit does not perform side-effecting work.  It
    does not invoke binaries, call APIs, or modify system state.  The Kit
    compiler is a pure function: YAML in, YAML out.

  \item[Kits are not plugins in a generic sense.] Gert does not have a
    ``plugin system.''  It has three distinct extension points---Domain Kits
    (authoring semantics), extensions (runtime capabilities), and tools
    (effectful actions)---each with a different contract, lifecycle, and trust
    model.  Calling any of these ``plugins'' obscures the architectural
    boundaries.

  \item[Kits are not templates or examples.] A Kit is not a collection of
    example runbooks.  It is a semantic layer: schemas, validators, a compiler,
    projections, and testing contracts.  An example runbook may demonstrate a
    Kit's vocabulary, but the Kit itself is the machinery that makes that
    vocabulary valid and executable.

  \item[Kits do not own execution.] The core runtime is the sole authority
    over step advancement, state management, governance enforcement, evidence
    capture, and trace writing.  A Kit may influence what the runtime executes
    (via lowering), but it never participates in how execution proceeds.

  \item[Kits do not replace the core schema.] The core \texttt{runbook/v2}
    schema remains the stable contract between authors, the parser, and the
    runtime.  Kit schemas are overlays that compile down to core; they do not
    bypass or replace it.

  \item[Kits are not cloud-native constructs.] Consistent with gert's
    local-first philosophy, Kits do not require network access, remote
    registries, or cloud services.  A Kit is a local artifact---a directory
    with a manifest, a compiler, and schemas---that runs wherever gert runs.
\end{description}
